\section{Schutz fremder und eigener Inhalte}

\subsection{Einführung und Beispiele}

\begin{example}
    \textbf{Aktuelle Fragestellungen und Beispiele:}
    \begin{itemize}
        \item Affen-Selfie (Wem gehören die Rechte?)
        \item KI-Training mit geschützten Daten
        \item KI ersetzt Online-Informationsdienste mit deren eigenen Inhalten
        \item Gibt es ein Opt-Out-Recht für KI-Training?
    \end{itemize}
\end{example}

\begin{remark}
    Verschiedene Gesetze regeln den Schutz von Inhalten und weisen unterschiedliche Ansätze auf. Relevante Rechtsgebiete sind unter anderem:
    \begin{itemize}
        \item Urheberrechtsgesetz
        \item Gesetz gegen den unlauteren Wettbewerb
        \item Obligationenrecht
        \item Markenschutz
        \item Datenschutz
    \end{itemize}
\end{remark}

\subsection{Das Urheberrecht}

\begin{definition}
    \textbf{Urheberrecht:} Das Urheberrecht schützt eine schöpferische Leistung. Der Urheber hat das Recht, über die Verwendung dieser Leistung zu bestimmen.
\end{definition}

\begin{remark}
    \textbf{Was schützt das Urheberrecht?} \\
    Der Weg von der Idee zur Verkörperung: Idee / Gedanke / Information $\rightarrow$ Formulierung / Umsetzung $\rightarrow$ Verkörperlichung. \\
    Das Urheberrecht schützt die \textbf{Formulierung bzw. Umsetzung}. \\
    \textit{"Schutzfähig ist nur der bestimmte Ausdruck eines Gedankens. Erst seine konkrete Verkörperung macht das Werk aus."}
\end{remark}

\begin{definition}
    \textbf{Bedingungen für urheberrechtlichen Schutz:}
    Damit ein Werk geschützt ist, müssen grundsätzlich folgende Bedingungen erfüllt sein:
    \begin{itemize}
        \item \textbf{Geistige Schöpfung:} Es muss ein Gestaltungswille bzw. eine bewusste Kreation vorliegen.
        \item \textbf{Individueller Charakter:} Das Werk muss sich abheben und kreativ sein.
    \end{itemize}
\end{definition}

\begin{remark}
    \textbf{Neuerung 2020 (Lichtbildschutz):}
    Gemäss dem neuen Gesetz von 2020 benötigen Fotografien (z.B. von dreidimensionalen Objekten) keinen individuellen Charakter mehr, um rechtlich geschützt zu sein.
\end{remark}
